\makeabstract
    {
    Using Single-Channel Electrophysiology to Identify Peptide Aggregation
    }
    {
    Cara Wagner-Oke\textsuperscript{1}, Devika Vikrama\textsuperscript{2}, Liviu Movileanu\textsuperscript{2}
    }
    {
    \textsuperscript{1} Biochemistry \& Biophysics Program, Amherst College, Amherst, MA\\
\textsuperscript{2} Department of Physics, Syracuse University, Syracuse, NY
    }
    {
    Nanopore sensors are a powerful tool for identifying and characterizing biomarkers, which can be used for disease detection. One of the most common causes of amyotrophic lateral sclerosis and frontotemporal dementia is an expansion of the hexanucleic repeat GGGGCC, which results in the production of dipeptide repeat proteins (DPRs). This overexpression results in the toxic aggregation of DPRs. Here we are working to develop a nanopore sensor for the detection of proline-arginine repeats (poly-PR), which is one of the most toxic DPR products.
    }
    {
    We expressed the nanopore z-tFhuA and inserted it into a synthetic lipid bilayer. We added poly-PR, which upon translocation of the nanopore, produces a distinct electrical signal. We recorded signal frequency and dwell time at different concentrations of poly-PR. By analyzing nanopore electrical recordings, we examined how frequency and dwell time relate to poly-PR concentration and aggregation. 
    }
    {
    We found that an increase in poly-PR concentration increases the nanopore signal frequency and that there are three distinct interaction lengths, which may correspond to states of poly-PR aggregation.
    }
    {
    We successfully created a nanopore sensor that interacts with poly-PR to produce a distinct electrical signal. These electrical signals provide data at the single-molecule level of nanopore-poly-PR interactions, allowing us to develop a technique for identifying early-stage onset of amyotrophic lateral sclerosis and frontotemporal dementia.
    }

    \newpage
    